% cap3.tex
\chapter{Metodología}
\label{ch:metodo}

{\color{blue} En esta sección se debe incluir el diseño experimental y la metodología general utilizados, junto con su explicación y el detalle de materiales, equipos y métodos teóricos y experimentales. Máximo 15 páginas.}

\section{Ejemplo de texto}
En esta sección describimos el proceso de recopilación de los datos necesarios para la investigación, que incluyó corpus jurídicos chilenos e internacionales, además de proyectos del Servicio de Evaluación Ambiental de Chile.

Para cada modelo se crea una estructura JSON que almacena las métricas de cada época durante el ajuste fino. Esto facilita el seguimiento y análisis del rendimiento en distintos conjuntos de datos. Cada modelo y conjunto de datos tiene un archivo JSON propio. La estructura registra, entre otras, las siguientes métricas:

\begin{figure}[H]
\centering
\includegraphics[scale=0.5]{images/JSON_ft.png}
\caption[Ejemplo de JSON creado durante el ajuste fino]{Ejemplo de JSON creado durante el ajuste fino. Fuente: elaboración propia}
\label{fig:JSON Fine Tune Stago}
\end{figure}

\begin{itemize}
   \item 10: número de la época de ajuste fino.
   \item tLoss: pérdida de entrenamiento observada en la última época.
   \item eLoss: pérdida de evaluación sobre el conjunto de validación.
   \item Acc/Recall micro: exactitud y exhaustividad calculadas mediante promedio micro.
   \item Recall macro: promedio de la exhaustividad de todas las clases.
   \item F1 macro: promedio de la medida F1 de todas las clases.
   \item Train Time: duración del entrenamiento.
   \item Eval Time: duración de la evaluación.
   \item Confusion Matrix (CFM): detalle de verdaderos positivos, falsos positivos, verdaderos negativos y falsos negativos por clase.
\end{itemize}

El formato JSON estructurado facilita la comparación entre épocas, modelos y conjuntos de datos. La matriz de confusión también se representa mediante un mapa de calor para mejorar su interpretación.

\begin{figure}[H]
\centering
\includegraphics[scale=0.8]{images/Example_CM.png}
\caption[Ejemplo de matriz de confusión]{Ejemplo de matriz de confusión. Fuente: elaboración propia}
\label{fig:Confusion Matrix Example}
\end{figure}
